%
\vfill
\begin{quote}
  These lecture notes bridge a gap between introductory quantum field theory (QFT) courses and state-of-the-art research in scattering amplitudes. They cover the path from basic definitions of QFT to amplitudes relevant for processes in the Standard Model of particle physics. The book begins with a concise yet self-contained introduction into QFT, including perturbative quantum gravity. It then presents modern methods for calculating scattering amplitudes, focusing on tree-level amplitudes, loop-level integrands and loop-integration techniques. These methods help reveal intriguing relations between gauge and gravity amplitudes, and are of increasing importance for obtaining high-precision predictions for collider experiments, such as those at CERN's Large Hadron Collider, as well as for foundational mathematical physics studies in QFT, including recent applications to gravitational wave physics.
\vskip 0.2cm
These course-tested lecture notes include numerous exercises with detailed solutions. Requiring only minimal knowledge of QFT, they are well-suited for MSc and PhD students as a preparation for research projects in theoretical particle physics. They can be used as a one-semester graduate level course, or as a self-study guide for researchers interested in fundamental aspects of QFT. Supplementary material, \textsc{Mathematica} notebooks, corrections and further information are provided and maintained at the dedicated website
%
\begin{center}
\href{https://scattering-amplitudes.mpp.mpg.de/scattering-amplitudes-in-qft/}{https://scattering-amplitudes.mpp.mpg.de/scattering-amplitudes-in-qft/}\,.
\end{center}
\end{quote}
